\documentclass[conference]{IEEEtran}
\IEEEoverridecommandlockouts

\usepackage{cite}
\usepackage{amsmath,amssymb,amsfonts}
\usepackage{graphicx}
\usepackage{textcomp}
\usepackage{booktabs}
\usepackage{array}
\usepackage{url}

\begin{document}

\title{Passive RF-Silent Drone Detection via Neuromorphic Vision and
Acoustic Fusion of the Propeller Blade-Pass Frequency}

\author{\IEEEauthorblockN{Giray Yillikci}
\IEEEauthorblockA{\textit{Independent Researcher} \\
gyillikci@gmail.com}
\and
\IEEEauthorblockN{\.{I}sa \"{O}zy{\i}ld{\i}r{\i}m}
\IEEEauthorblockA{\textit{Independent Researcher}}}

\maketitle

\begin{abstract}
We investigate passive, RF-silent drone detection by exploiting the fact
that a spinning propeller emits a characteristic blade-pass frequency
(BPF) simultaneously as optical flicker, sensed by a neuromorphic event
camera, and as acoustic pressure, sensed by a microphone array. To our
knowledge no prior work fuses these two modalities for counter-UAS. We
present a working prototype that extracts a per-pixel flicker frequency
from a Prophesee EVK4-HD event camera, isolates propeller regions by
connected-component clustering and temporal tracking, extracts the
acoustic BPF and a direction-of-arrival estimate from a four-microphone
array, and fuses the two with a sequential Bayesian updater using
harmonic-aware frequency cross-validation. The central finding is
methodological: detection \emph{architecture}, not parameter tuning,
determines success. A grid-FFT baseline failed at range regardless of
tuning, whereas per-pixel frequency extraction with connected-component
clustering gave stable, zero-false-positive detection with first-detect
latency reduced from 600--1200\,ms to ${\sim}100$\,ms after real-time
optimization. Cross-modal frequency agreement worked reliably; cross-modal
\emph{spatial} calibration did not converge indoors ($R^2 = 0.14$), which
we report as an instructive negative result.
\end{abstract}

\begin{IEEEkeywords}
event camera, neuromorphic vision, drone detection, sensor fusion,
blade-pass frequency, Bayesian inference
\end{IEEEkeywords}

\section{Introduction}
Unauthorized consumer drones are now an operational threat: the 2018
Gatwick incident diverted ${\sim}1000$ flights, and drones flying
pre-programmed GPS missions emit no radio signal, making them invisible to
RF scanning---the dominant commercial technology. Since only a few US
federal agencies may lawfully jam a drone, most operators are restricted to
\emph{detection}, motivating passive sensing. Every existing modality has a
structural blind spot (Table~\ref{tab:modalities}). Prior counter-UAS work
spans RF fingerprinting \cite{guvenc2018,ezuma2020} and acoustic signature
detection with array localization \cite{mezei2015,shi2018}; we are not
aware of any that fuses an event camera with a microphone array.

\begin{table}[t]
\caption{Structural Weaknesses of Existing Detection Modalities}
\label{tab:modalities}
\centering
\small
\begin{tabular}{@{}lp{5.4cm}@{}}
\toprule
\textbf{Modality} & \textbf{Structural weakness} \\
\midrule
RF scanning & Blind to autonomous / RF-silent drones \\
Radar & Confuses small drones with birds; costly \\
RGB / thermal & Motion blur, ${\sim}60$\,dB range, low-light/fog failure \\
Acoustic only & Short range; 10--20\,\% false positives in noise \\
\bottomrule
\end{tabular}
\end{table}

\textbf{Hypothesis.} The BPF appears in both the optical and acoustic
domains; requiring \emph{cross-modal frequency agreement} suppresses
false-alarm sources that live in only one modality (flickering luminaires
are silent; HVAC does not flicker), driving the joint false-positive
probability down.

\textbf{Contributions.} (i) A prototype---to our knowledge the
first---fusing an event camera with a microphone array for drone detection
via the shared BPF. (ii) An architectural result: per-pixel frequency plus
connected-component clustering succeeds where grid-FFT fails at range
regardless of tuning. (iii) Harmonic-aware cross-modal frequency validation
and gated Bayesian fusion. (iv) Real-time engineering cutting first-detect
latency $6$--$12\times$. (v) An honest account of negative results.

\section{Physical Principle}
An event camera reports asynchronous per-pixel brightness-change events
$(x,y,t,\text{polarity})$ with microsecond timing and $>$120\,dB dynamic
range \cite{lichtsteiner2008,gallego2020,finateu2020}; a 30\,fps camera samples a
200\,Hz BPF below the Nyquist rate and blurs the blades, whereas an event
camera measures the per-pixel oscillation directly. A propeller with $B$
blades at shaft rate $f_s$ modulates both light and sound at
\begin{equation}
f_{\mathrm{BPF}} = B f_s = B\,\mathrm{RPM}/60,
\label{eq:bpf}
\end{equation}
typically 80--300\,Hz for common multirotors (1--2.5\,kHz for racing
quads). The signature is unavoidable, narrowband, and rare in nature, so
requiring the two channels to agree (up to harmonic ambiguity,
Section~\ref{sec:fusion}) rejects single-modality noise.

\section{Method}

\subsection{Visual Pipeline}
\textbf{Grid-FFT baseline (failed).} Dividing the $1280\times720$ sensor
into a $16\times16$ grid, counting events per cell, and running a per-cell
FFT detected the propeller only at centimeter range. At distance a
propeller's few pixels are diluted by hundreds of noise pixels in the same
cell; relaxing thresholds merely promoted noise to 25--43 false
``regions'' per frame. \emph{Tuning cannot fix a wrong architecture.}

\textbf{Per-pixel frequency $+$ connected components.} We instead threshold
the SDK per-pixel frequency map to a binary mask, dilate, run
connected-component labeling, and keep clusters of $\ge 2$ pixels whose
frequency coefficient of variation is $\le 0.3$. Clustering adapts to the
propeller's true shape---even 2--3 consistent pixels are meaningful---so
there is no signal dilution. A single-pixel minimum produced 13--17 false
detections per frame, confirming that spatial coherence ($\ge 2$ connected
pixels) is the primary false-positive filter. A temporal tracker with
greedy association (spatial $<$100\,px, frequency $<$30\,\%) and
velocity-predictive matching confirms only persistent tracks. On the bench
propeller this gave one stable track at 270\,Hz BPF (5400\,RPM, 3 blades),
locked position, and zero false positives.

\textbf{Real-time optimization.} Full-resolution morphology and connected
components at 25\,Hz cost ${\sim}315$\,ms/s of CPU, causing growing lag. A
$4\times$ downscale before morphology (statistics still on full-resolution
regions of interest), in-place drawing, and early-exit on empty masks
removed all drift. A production pass then pre-allocated all buffers
(eliminating garbage-collection jitter) and replaced hard hit-count
confirmation with quality-weighted Bayesian confidence: each observation
contributes $p = 0.65 + q(0.85-0.65)$, with $q$ from pixel count and
frequency consistency, accumulated as $P = 1-(1-p)^n$, decayed when idle,
confirming at $P \ge 0.95$. Strong detections then confirm in 2 frames and
weak ones in 4, preserving zero false positives.

\subsection{Acoustic Pipeline}
A rolling 0.2\,s window per channel is Hann-windowed and transformed; the
four magnitude spectra are incoherently averaged (${\sim}6$\,dB gain) and
the peak picked in the target band with an SNR gate. Direction of arrival
uses frequency-weighted GCC-PHAT \cite{knapp1976} per opposing 64\,mm mic
pair with $4\times$ zero-padded sub-sample delay, parabolic interpolation,
a peak/second-peak confidence gate, and a median filter, giving usable
azimuth on steady tones (${\sim}\pm 10^{\circ}$; the planar array cannot
resolve elevation). ODAS \cite{grondin2019} and the array's speech-tuned
firmware DOA both failed on impulsive propeller acoustics, motivating the
custom estimator.

\subsection{Fusion}
\label{sec:fusion}
One sensor often reports a harmonic of the other, so we search pairs
$(n_v,n_a)$, $1\le n_v,n_a\le 4$, and declare a match when
$|f_v/n_v - f_a/n_a|/(f_a/n_a) < 0.05$ with confidence $1/(n_v n_a)$. A
sequential Bayesian detector starts at prior $P=0.001$, updates per
modality with sigmoid likelihoods, adds a frequency-match bonus, and decays
toward the prior absent evidence; detection fires at $P>0.8$. Quality,
temporal-association ($\le 200$\,ms), and angle-residual ($\le 35^{\circ}$)
gates prevent overconfident fusion. In a live run with no propeller in view
the detector latched on strong acoustic tones via the audio-only path and
decayed when they stopped---single-modality dominance is a real failure
mode the gates must police.

\section{Results}
The headline result is latency (Table~\ref{tab:latency}): first-detect
falls from 600--1200\,ms to ${\sim}100$\,ms with jitter reduced
${\sim}5\times$. The final configuration tracked multi-pixel clusters over
$80{+}$ consecutive hits with zero false positives on ${\sim}30$\,s
recordings, and cross-modal frequency agreement fired reliably.

\begin{table}[t]
\caption{Visual Detector Performance, v2 vs.\ v3}
\label{tab:latency}
\centering
\begin{tabular}{@{}lccc@{}}
\toprule
\textbf{Metric} & \textbf{v2} & \textbf{v3} & \textbf{Gain} \\
\midrule
First detect (strong) & 600--1200\,ms & ${\sim}100$\,ms & $6$--$12\times$ \\
First detect (weak) & 600--1200\,ms & ${\sim}200$\,ms & $3$--$6\times$ \\
Steady-state update & 100--140\,ms & 50--90\,ms & ${\sim}2\times$ \\
Per-frame analysis & ${\sim}1.5$\,ms & ${\sim}0.8$\,ms & ${\sim}2\times$ \\
GC jitter & $\pm 5$\,ms & $<\pm 1$\,ms & ${\sim}5\times$ \\
\bottomrule
\end{tabular}
\end{table}

\textbf{Azimuth calibration (negative result).} Fitting a linear
visual$\rightarrow$acoustic azimuth map by sweeping a co-located source
yielded $R^2 = 0.14$ (RMSE $0.70^{\circ}$): the acoustic DOA spanned only
${\sim}\pm 1.2^{\circ}$ against a ${\sim}\pm 25^{\circ}$ visual sweep,
consistent with reverberation or an insufficient sweep. Cross-modal
\emph{spatial} fusion is therefore unvalidated; cross-modal \emph{frequency}
fusion is validated. Projected ranges (${\sim}50$--150\,m,
acoustic-limited) and the sub-0.1\,\% false-positive target are
design-derived, pending outdoor validation.

\section{Discussion and Limitations}
Three lessons generalize. \emph{Architecture beats tuning}: the decisive
change was replacing grid-FFT counting with per-pixel frequency plus
clustering. \emph{Spatial coherence is the false-positive filter}: two
connected consistent pixels are meaningful, one is not.
\emph{Probabilistic beats deterministic confirmation}: quality-weighted
Bayesian confidence confirms strong detections $3$--$6\times$ faster than a
hard hit count at the same false-positive rate.

Limitations: all detection results are indoor, single-source,
${\sim}30$\,s recordings; the planar array cannot resolve elevation;
harmonic ambiguity (2-blade at 100\,Hz $\equiv$ 4-blade at 50\,Hz) needs
external metadata; and there is no hardware inter-sensor clock
(${\sim}100$\,ms association uncertainty). Planned work: a headless C++
port for deterministic sub-millisecond timing, a non-planar array for
elevation, hardware synchronization, multi-target tracking, and outdoor
calibration with a wider sweep. The validated contributions stand
independently of the range and false-positive projections, which await
outdoor validation on labeled data.

\begin{thebibliography}{00}
\bibitem{lichtsteiner2008} P. Lichtsteiner, C. Posch, and T. Delbruck,
``A $128\times128$ 120 dB 15 $\mu$s latency asynchronous temporal contrast
vision sensor,'' \emph{IEEE J. Solid-State Circuits}, vol. 43, no. 2,
pp. 566--576, 2008.
\bibitem{gallego2020} G. Gallego \emph{et al.}, ``Event-based vision: A
survey,'' \emph{IEEE Trans. Pattern Anal. Mach. Intell.}, vol. 44, no. 1,
pp. 154--180, 2022.
\bibitem{finateu2020} T. Finateu \emph{et al.}, ``A $1280\times720$
back-illuminated stacked temporal contrast event-based vision sensor,'' in
\emph{IEEE ISSCC}, 2020, pp. 112--114.
\bibitem{guvenc2018} I. Guvenc \emph{et al.}, ``Detection, tracking, and
interdiction for amateur drones,'' \emph{IEEE Commun. Mag.}, vol. 56,
no. 4, pp. 75--81, 2018.
\bibitem{ezuma2020} M. Ezuma \emph{et al.}, ``Micro-UAV detection and
classification from RF fingerprints using machine learning,'' \emph{IEEE
Access}, vol. 8, pp. 7654--7668, 2020.
\bibitem{mezei2015} J. Mezei, V. Fiaska, and A. Moln\'ar, ``Drone sound
detection,'' in \emph{IEEE CogInfoCom}, 2015, pp. 333--338.
\bibitem{shi2018} Z. Shi \emph{et al.}, ``An acoustic-based surveillance
system for amateur drones detection and localization,'' \emph{IEEE Trans.
Veh. Technol.}, vol. 69, no. 3, pp. 2731--2739, 2020.
\bibitem{knapp1976} C. Knapp and G. Carter, ``The generalized correlation
method for estimation of time delay,'' \emph{IEEE Trans. Acoust., Speech,
Signal Process.}, vol. 24, no. 4, pp. 320--327, 1976.
\bibitem{grondin2019} F. Grondin and F. Michaud, ``Lightweight and
optimized sound source localization and tracking methods,'' \emph{Robot.
Auton. Syst.}, vol. 113, pp. 63--80, 2019.
\end{thebibliography}

\end{document}
